\documentclass[11pt]{article}
\usepackage{amsmath,amssymb,amsthm}
\newtheorem{proposition}{Proposition}
\title{Exact Moment and Fourier Identities for Oriented Fractions}
\author{}
\date{November 24, 2026}
\begin{document}
\maketitle

\paragraph{Status and scope.}
These are complete written identities, checked on exact rational
examples by \texttt{positive-cube-moments-check.py}. They are not Lean
theorems and do not establish either requested asymptotic bound.
The Fourier resolvent statement applies to every oriented positive
affine denominator. The odd resolvent statement applies to the full
odd-target normalization already checked in
\texttt{HeadComplexity.OddKernel.of\_fracComputable}.
No bounded margin, bounded numerator, or separation between heads is
assumed. The last paragraph identifies the cancellation gap explicitly.

For a function on $\{0,1\}^n$, write $[f]_S$ for its coefficient of
$x_S$ in its unique multilinear extension. Let
\[
 B(x)=b_0+\sum_i b_ix_i,\qquad
 A(x)=a_0+\sum_i a_ix_i,\qquad b_0,b_i>0.
\]

\begin{proposition}[Positive moment formula]
Put
\[
 m_S=\int_0^\infty e^{-b_0t}
                \prod_{i\in S}(1-e^{-b_it})\,dt.
\]
Then $[1/B]_S=(-1)^{|S|}m_S$. For $S\ne\varnothing$, put
\[
 I_{S,i}=\int_0^\infty e^{-(b_0+b_i)t}
             \prod_{k\in S\setminus\{i\}}(1-e^{-b_kt})\,dt>0.
\]
The exact affine-numerator formula is
\[
 (-1)^{|S|}[A/B]_S
       =\sum_{i\in S}(a_0b_i/b_0-a_i)I_{S,i}.                 \tag{1}
\]
\end{proposition}
\begin{proof}
Use $1/B(x)=\int_0^\infty e^{-b_0t}\prod_i e^{-b_ix_it}\,dt$ and
$e^{-b_ix_it}=1-(1-e^{-b_it})x_i$ at Boolean inputs. All sums are finite
and all integrals converge. Multiplication by $x_i$ and the identity
$x_i^2=x_i$ show that the coefficient on the left of (1) is
\[
 a_0m_S+\sum_{i\in S}a_i(m_S-m_{S\setminus\{i\}})
             =a_0m_S-\sum_{i\in S}a_iI_{S,i}.
\]
Integration of the derivative of
$e^{-b_0t}\prod_{i\in S}(1-e^{-b_it})$ gives
$b_0m_S=\sum_{i\in S}b_iI_{S,i}$. Its boundary values are zero at
both endpoints because $S$ is nonempty. Substitution proves (1).
\end{proof}

\paragraph{Both orientations and constant denominators.}
A strictly negative-slope denominator positive on the cube becomes
the displayed positive form under $y_i=1-x_i$ for every coordinate.
Formula (1) then concerns coefficients in the $y$ basis; it is not
permissible to combine the two orientations as though their original
$x$ coefficients had the same alternating sign. A constant denominator
has only affine coefficients and is handled separately. Thus the
identity respects all original denominator cases.

\begin{proposition}[A consequence and its limitation]
The positive tensor $m$ is log-supermodular:
\[
             m_Sm_T\leq m_{S\cap T}m_{S\cup T}.              \tag{2}
\]
Nevertheless positivity and (2) alone impose no restriction on the
strict signs of a difference of two such positive tensors: for every
array $\sigma_S\in\{-1,1\}$ there are strictly positive
log-supermodular arrays $u,v$ with
$\operatorname{sgn}(u_S-v_S)=\sigma_S$ for every $S$.
The latter arrays are not asserted to be reciprocal moment tensors.
\end{proposition}
\begin{proof}
For (2), normalize the finite measure
$e^{-b_0t}\prod_{i\in S\cap T}(1-e^{-b_it})\,dt$. The two remaining
products for $S\setminus T$ and $T\setminus S$ are increasing
functions of $t$. Their covariance is nonnegative, since it is half
the integral over two independent samples of the product of their
two differences. Undoing the normalization gives (2).

For the second assertion take $v_S=2^{|S|^2}$. If $S,T$ are
incomparable, then
\[
 |S\cup T|^2+|S\cap T|^2-|S|^2-|T|^2
                 =2|S\setminus T|\,|T\setminus S|>0.
\]
All corresponding inequalities are therefore strict. There are only
finitely many. For sufficiently small $\epsilon>0$, the array
$u_S=v_S+\epsilon\sigma_S$ remains positive and satisfies all of
them. For comparable sets the inequality is an identity for every
array. Finally $u_S-v_S=\epsilon\sigma_S$.
\end{proof}

\paragraph{Fourier notation.}
On $\{-1,1\}^n$ put $\chi_S(t)=\prod_{i\in S}t_i$, and let
$\widehat f(S)=2^{-n}\sum_t f(t)\chi_S(t)$. Let $T_i$ act on
vectors indexed by subsets by $(T_ic)_S=c_{S\mathbin\triangle\{i\}}$.
These commuting matrices satisfy $T_i^2=I$.

\begin{proposition}[Positive Fourier resolvent]
Suppose $w_i\geq0$ and $\tau=\sum_iw_i<1$. Put
$L_w=\sum_iw_iT_i$. For
\[
 R(t)=\frac{a_0+a\cdot t}{1-w\cdot t}
\]
the Fourier vector is
\[
 \widehat R=(I-L_w)^{-1}
           \left(a_0e_\varnothing+\sum_i a_ie_{\{i\}}\right),
 \qquad (I-L_w)^{-1}=\sum_{k\geq0}L_w^k.                   \tag{3}
\]
Every matrix entry of the inverse is nonnegative. Every original
uniformly oriented positive affine denominator has this form after
positive scalar normalization and possibly complementing all inputs.
\end{proposition}
\begin{proof}
Multiplication by $t_i$ toggles Fourier index $i$. Applying the
Fourier transform to $(1-w\cdot t)R=a_0+a\cdot t$ gives (3).
The row-sum norm of $L_w$ is $\tau<1$, so the geometric matrix series
converges and is the inverse. Its entries are nonnegative.
For a general centered denominator $d+v\cdot t$, positivity at all
vertices gives $d>\sum_i|v_i|$. Uniform slope orientation permits a
single global complement to make all its slopes nonpositive; then
divide by $d$. Zero slopes, including all slopes zero, cause no issue.
\end{proof}

\begin{proposition}[Odd coefficients are a killed-walk extension]
Let
\[
 G(t)=\frac{u\cdot t}{1-(w\cdot t)^2},\qquad
 w_i\geq0,\qquad \tau=\sum_iw_i<1.
\]
Its Fourier vector on odd subsets satisfies
\[
 c=(I-L_w^2)^{-1}\sum_i u_ie_{\{i\}}.                      \tag{4}
\]
Let $D$ be the singletons and $E$ the odd subsets of size at least
three. With $M=L_w^2$ restricted to odd subsets,
\[
 c_E=P(w)c_D,\qquad
 P(w)=(I-M_{EE})^{-1}M_{ED}\geq0.                          \tag{5}
\]
For each $S\in E$ its row mass obeys
\[
 \sum_{i\in D}P(w)_{S,i}\leq\tau^{|S|-1},\qquad
 |c_S|\leq\tau^{|S|-1}\max_i|c_{\{i\}}|.                 \tag{6}
\]
\end{proposition}
\begin{proof}
Equation (4) follows by Fourier transformation just as in (3).
The matrix $M$ is nonnegative with every row sum $\tau^2<1$.
Its $E$ equations have zero source, giving (5). The inverse in (5)
is again a convergent series of nonnegative matrices.

If $\tau=0$, $M=0$ and the claims follow. Otherwise put
$v_S=\tau^{|S|-1}$ on odd subsets. Then $v_D=1$.
Every nonzero entry $M_{S,T}$ has $|T|\geq|S|-2$, so
$v_T\leq\tau^{-2}v_S$. The row-sum identity implies
$(Mv)_S\leq v_S$. Restricting this inequality to $E$ gives
$(I-M_{EE})v_E\geq M_{ED}1$. Multiply by its nonnegative inverse
to obtain $v_E\geq P(w)1$, which is the first assertion in (6).
The second follows from (5) and the triangle inequality.
\end{proof}

\paragraph{What remains missing for the many-head lower bound.}
For a sum of odd kernels, (6) bounds a high coefficient by
\[
 \sum_h \tau_h^{|S|-1}\max_i|\widehat G_h(\{i\})|.
\]
There is no bound here on the individual first-level coefficients in
terms of the first-level coefficients of the sum, its head count,
or its strict signs. Different positive extension matrices in (5)
permit their first-level sources to cancel while higher coefficients
remain. The exact two-kernel three-variable parity example in
\texttt{odd-positive-kernel-normalization.tex} already exhibits this.
Consequently neither truncating (3) nor applying (6) with a bound
on the total score would be justified without an additional,
cancellation-safe argument. No such argument is supplied here.

\begin{proposition}[Promised shattering survives odd normalization]
Suppose a sum of $H\geq1$ positive-slope admissible affine ratios
has specified strict signs on a finite subset $S\subseteq\{0,1\}^n$.
There is an antipodally odd sum of $2H$ admissible ratios in $n+1$
variables with the same strict signs on $\{0\}\times S$.
Here odd means $G(1-u,1-x)=-G(u,x)$.
\end{proposition}
\begin{proof}
Absorb any scalar bias into the first affine numerator, and write
$F=\sum_h A_h/B_h$. For $T>0$ put
\[
 H_T(u,x)=\sum_h\frac{A_h(x)}{B_h(x)+Tu},\qquad
 G_T(u,x)=H_T(u,x)-H_T(1-u,1-x).
\]
The first $H$ denominators have strictly positive slopes, and the
last $H$ have strictly negative slopes. All are positive on the full
cube. Each numerator remains affine. The displayed odd identity is
exact. On $u=0$,
$G_T(0,x)=F(x)-H_T(1,1-x)$, and the second term tends uniformly to
zero on the finite cube as $T\to\infty$. Choose $T$ large enough to
preserve the finite positive margin on $S$.
\end{proof}

In particular the forward majority-order grid compiler already
audited in \texttt{exponent-mechanisms.tex} gives promised-domain odd
scores with $\Omega(n^2)$ shattered points and $O(n/\log n)$ heads.
A generic $O(nH)$ VC bound for odd scores would therefore be false
as well. A lower bound specifically for the entire prescribed cube
is not contradicted by this construction.
\end{document}
